\subsection{一个 Hochschild 上同调计算}\label{sec:20d}
我们暂不直接对 $A$ 上的 $\Ai$ 结构进行分类，而先转向一个稍大的代数 $\overline A$。它是扩充箭图——比原来多一个顶点——
\SourceEquationDiagram{diagram-20-10}{0.415384168}{20.10}
模去关系
\[
\begin{gathered}
(i+2,i+1,i)=0,\qquad(i-2,i-1,i)=0,\\
(i,i+1,i)+(i,i-1,i)=0,\qquad(0,1,0)=0
\end{gathered}
\]
所得的商。由构造，$\overline A$ 是 $\overline R=\K^{m+1}$ 上的线性代数。或许令人意外的是，这个扩充代数比 $A$ 具有更好的性质。它的上同调维数有限，而且还是 Koszul 代数。下面解释后一性质的含义。

由于 $\overline A$ 是路径代数模去一个关于次数和路径长度都齐次的理想所得的商，它同样具有双分次。因此，Ext 代数
\begin{equation}\label{eq:20.11}
\operatorname{Ext}_{\overline A}(\overline R,\overline R)
\end{equation}
\SourcePage{317}
先验地带有三重分次：上同调次数、由 $\overline A$ 中元素的次数导出的内部分次，以及同样由 $\overline A$ 中相应分次导出的长度分次。考察计算\eref{20.11}所用的约化条形构造，显然，上链的上同调分次总不超过其长度分次。此外，可以显式描述上链复形中这两个分次相等的部分。结果如下。令 $\overline A^*$ 为 $\overline A$ 的二次对偶；其生成元 $(i,i\pm1)^*$ 与 $(i\pm1,i)$ 对偶，次数为
\[
|(i,i\pm1)^*|=1-|(i\pm1,i)|,
\]
并满足关系
\[
(i,i+1,i)^*+(i,i-1,i)^*=0,\qquad(m,m-1,m)^*=0.
\]
于是，将条形上链限制到生成元 $(i\pm1,i)$ 上，便定义了一个典范同态
\begin{equation}\label{eq:20.12}
\operatorname{Ext}_{\overline A}(\overline R,\overline R)\longrightarrow\overline A^*.
\end{equation}
在本例中，由\bcite{85}，命题~2.1 给出的单模的投射分解可知，Ext 代数的三重分次部分地集中在对角线上，即所有非零分次部分的上同调分次都等于长度分次。这恰好是\bcite{110}意义下的 Koszul 代数性质；这里稍作推广，以半单环 $\overline R$ 而非域作为基环。该性质等价于\eref{20.12}为同构。\znote{309}{“三重分次部分地对角”指：Ext 代数中每个非零分次部分都满足上同调次数等于长度分次。这与上文“上链的上同调次数总不超过长度分次”合起来，说明整个 Ext 代数由对角部分决定，于是 (20.12) 是同构就等价于 Koszul 性。此处基环取半单环 $\overline R=\K^{m+1}$ 而非域，所以作者说这是\bcite{110}中定义的轻微推广。[来源：A，原书本节与所引文献。]}

接着，考虑 $\overline A$ 的约化 Hochschild 上链复形，它由从 $\overline A/\overline R$ 到 $\overline A$ 的多线性映射组成。可按映射的目标中的路径长度过滤这个复形。所得谱序列的 $E_1$ 页可写成
\begin{equation}\label{eq:20.13}
\operatorname{Ext}_{\overline A}(\overline R,\overline R)
\otimes_{\overline R\otimes\overline R}\overline A
\cong\overline A\otimes_{\overline R\otimes\overline R}\overline A^*.
\end{equation}
由于 $\overline A$ 及其对偶均为有限维，这个谱序列显然收敛到 $\HH(\overline A,\overline A)$。从图形上看，$E_1$ 页的生成元可以视为由两段组成的闭路径，两段分别位于箭图\eref{20.10}及其对偶中。具体地，它们可写成
\begin{equation}\label{eq:20.14}
(i_l,\ldots,i_0)\otimes(j_0,\ldots,j_r)^*,\qquad j_r=i_l,\quad j_0=i_0.
\end{equation}
分次问题可以如下理清。假设我们关心 $\HH^{r+s}(\overline A,\overline A)^s$，它是由内部次数为 $s$ 的 $r$ 线性 Hochschild 上闭链表示的群；这里的记号与\secref{1f}相同。那么，\eref{20.13}中对它有贡献的项，是\eref{20.14}中具有给定 $r$ 并满足
\[
|(j_0,\ldots,j_r)^*|+|(i_l,\ldots,i_0)|=r+s
\]
的那些表达式。由箭图的分次方式，得到 $l=2s+r$。
\StatementNumber{2}\begin{lemma}\label{st:20.2}
当 $r\ge5$ 或 $r=3$ 时，$\HH^2(\overline A,\overline A)^{2-r}$ 为零；当 $r=4$ 时，它为 $\K$。
\end{lemma}
\SourcePage{318}
\begin{proof}[证明提要]
当 $r\ge5$ 时，结论直接成立，因为有贡献的项\eref{20.14}必须满足 $l=4-r<0$。$r=3,4$ 的情形略为复杂；这里略去细节，但记下一个稍后要用的推论。\eref{20.13}的边缘同态给出映射
\begin{equation}\label{eq:20.15}
\HH^2(\overline A,\overline A)^{-2}\longrightarrow
\overline R\otimes_{\overline R\otimes\overline R}(\overline A^*)^2.
\end{equation}
这里，目标空间由 $\overline A^*$ 中次数为二的元素组成，$\overline R$ 的幂等元在这些元素上对角地作用，即左、右作用相同。由计算可知，\eref{20.15}是单射，其像由下列元素的倍数组成；该元素属于 $\overline A^*$ 的中心：
\begin{equation}\label{eq:20.16}
\begin{aligned}
&(0,1,2,1,0)^*+(1,2,3,2,1)^*+\cdots\\
&\hspace{4em}+(m-2,m-1,m,m-1,m-2)^*.
\end{aligned}
\end{equation}
\end{proof}

结合一般分类理论，\stref{引理}{20.2}表明，在固定上同调的拟同构意义下，满足 $H(\overline{\A})\cong\overline A$ 的 $\Ai$ 范畴 $\overline{\A}$ 至多构成一个单参数族。将参数记为 $q\in\K$；$q=0$ 对应形式的 $\Ai$ 结构。单凭\stref{引理}{20.2}本身，尚不能证明对应于 $q\ne0$ 的非形式结构存在。实际上，不同的非零参数值可通过对各个分次部分重新缩放而联系起来，所以它们并无实质区别。不过，暂且假设这样的结构存在，那么可以安排它们的第一个非平凡高阶乘积为 $\mu^4$。此外，由\eref{20.16}可知，当 $m\ge5$ 时，这一乘积满足\eref{20.8}的类似等式，且 $q\ne0$。特别地，此时若删去额外顶点对应的对象，由\stref{引理}{20.1}，所得全 $\Ai$ 子范畴仍不是形式的。
